
\documentclass{.local/lib/classes/ndvr-vector-image-standalone}

\title[NDVR thesis-abstract]{%
    Логическая модель физического представления видео%
}

\textwidth=40cm

\begin{document}

    %% My Types
    \newcommand{\umldtp}[1]{\textsf{#1}}
    \newcommand{\umlctp}[1]{\textsl{$\star$ #1}}
    \newcommand{\umlfld}[1]{\textsf{#1}}


    %% Tabular Styling
    \newcolumntype{F}[1]{>{\raggedleft\arraybackslash}m{#1}}
    \newcolumntype{T}[1]{>{\raggedright\arraybackslash}m{#1}}

    \setlength\dashlinedash{0.2pt}
    \setlength\dashlinegap{2.5pt}
    \setlength\arrayrulewidth{0.2pt}

    \renewcommand{\arraystretch}{1.2}
    \setlength{\tabcolsep}{0.1em}

    %% Tikz Uml Styling
    \colorlet{commonClass}{yellow!10}
    \colorlet{commonNote}{green!10}
    \colorlet{commonDraw}{darkgray}

    \colorlet{emphClass}{yellow!20}
    \colorlet{emphDraw}{black}

    \tikzumlset{
        %font=\normalsize,
        fill class=commonClass,
        fill note=commonNote,
        draw=commonDraw,
    }

    \tikzstyle{tikzuml inherit style}=[
        color=black,
        very thick,
        {-{Stealth[inset=0pt,scale=1.5,fill=white,angle'=45]}}
    ]

    \tikzstyle{tikzuml aggregation style}=[
        color=black,
        very thick,
        {{Diamond[scale=2,fill=white]}-}
    ]

    \tikzstyle{tikzuml composition style}=[
        color=black,
        very thick,
        {{Diamond[scale=2,fill=gray!90]}-}
    ]

    \tikzstyle{tikzuml implements style}=[
        color=black,
        very thick,
        {-{Stealth[inset=0pt,scale=1.5,fill=white,angle'=45]}}, 
        dashed
    ]

    %% Main Part
    \sf\singlespacing
    \begin{tikzpicture}[
        very thick,
        arrow/.style = {very thick,{Triangle[open]}-}
    ]
        \umlsimpleinterface{Видео}
        \umlnote [right=8em of Видео, width=15em] {Видео}{%
            На верхнем уровне иерархии
            видео представляет собой контейнер.
        }

        \umlclass[
            below=2em of Видео,
            draw=emphDraw,
            fill=emphClass,
        ]{Контейнер}{
            \begin{tabular}{cF{6em}cT{6em}}
                $+$ & \umlfld{Путь} 
                    &:& \umldtp{Строка};\\
                \hdashline
                $+$  & \umlfld{Длительность}
                     &:& \umldtp{Число секунд};\\
                \hdashline
                $+$  & \umlfld{Объём}
                     &:& \umldtp{Число байт};\\
                \hdashline
                $+$  & \umlfld{Битовая частота}
                     &:& \umldtp{Число};\\
                \hdashline
                $+$ & \umlfld{Описание} 
                    &:& \umlctp{Мета};\\
                \hdashline
                $+$ & \umlfld{Формат} 
                    &:& \umlctp{Формат контейнера};\\
                \hdashline
                $+$ & \umlfld{Опции} 
                    &:& \umldtp{Список опций};\\
                \hdashline
                $+$ & \umlfld{Потоки} 
                    &:& \umlctp{Список потоков};\\
            \end{tabular}
        }

        \umlclass[
            below=2em of Контейнер,
        ]{Мета}{
            \begin{tabular}{cF{6em}cT{4em}}
                $+$ & \umlfld{Название} 
                    &:& \umldtp{Строка};\\
                \hdashline
                $+$ & \umlfld{Описание} 
                    &:& \umldtp{Строка};\\
                \hdashline
                $+$ & \umlfld{Авторы}   
                    &:& \umldtp{Список строк};\\
                \hdashline
                $+$ & \umlfld{Метки} 
                    &:& \umlctp{Список строк};\\
            \end{tabular}
        }

        \umlclass[
            below=2em of Мета,
        ]{Формат контейнера}{
            \begin{tabular}{cF{6em}cT{4em}}
                $+$ & \umlfld{Название} 
                    &:& \umldtp{Строка};\\
                \hdashline
                $+$ & \umlfld{Описание} 
                    &:& \umldtp{Строка};\\
                \hdashline
                $+$ & \umlfld{Опции} 
                    &:& \umldtp{Список опций};\\
            \end{tabular}
        }

        \umlclass[
            below=2em of Формат контейнера,
        ]{Опция контейнера}{
            \begin{tabular}{cF{6em}cT{4em}}
                $+$ & \umlfld{Имя} 
                    &:& \umldtp{Строка};\\
                \hdashline
                $+$ & \umlfld{Значение} 
                    &:& \umldtp{Строка};\\
            \end{tabular}
        }

        \umlclass[
            type=abstract,
            right=3em of Контейнер,
            draw=emphDraw,
            fill=emphClass,
        ]{Абстрактный поток}{
            \begin{tabular}{cF{7.5em}cT{8em}}
                $+$  & \umlfld{Номер}  
                     &:& \umldtp{Число};\\
                \hdashline
                $+$  & \umlfld{Название}
                     &:& \umldtp{Строка};\\
                \hdashline
                $+$  & \umlfld{Описание}
                     &:& \umldtp{Строка};\\
                \hdashline
                $+$  & \umlfld{Язык}
                     &:& \umldtp{Строка};\\
                \hdashline
                $+$ & \umlfld{Время начала} 
                    &:& \umlctp{Точка времени};\\
                \hdashline
                $+$ & \umlfld{Время конца} 
                    &:& \umlctp{Точка времени};\\
                \hdashline
                $+$  & \umlfld{Битовая частота}
                     &:& \umldtp{Число};\\
                \hdashline
                $\#$ & \umlfld{Пакеты} 
                     &:& \umlctp{Список пакетов};\\
            \end{tabular}
        }

        \umlclass[
            below right=2em and 1em of Абстрактный поток.south west,
        ]{Визуальный поток}{
            \begin{tabular}{cF{7em}cT{6.5em}}
                $+$ & \umlfld{Ширина} 
                    &:& \umldtp{Число};\\
                \hdashline
                $+$ & \umlfld{Высота} 
                    &:& \umldtp{Число};\\
                \hdashline
                $+$ & \umlfld{Кодированная ширина} 
                    &:& \umldtp{Число};\\
                \hdashline
                $+$ & \umlfld{Кодированная высота} 
                    &:& \umldtp{Число};\\
                \hdashline
                $+$ & \umlfld{Видимое отношение сторон (DAR)} 
                    &:& \umldtp{(Число, Число)};\\
                \hdashline
                $+$ & \umlfld{Хранимое отношение сторон (SAR)} 
                    &:& \umldtp{(Число, Число)};\\
                \hdashline
                $+$ & \umlfld{Частота кадров} 
                    &:& \umldtp{Число};\\
                \hdashline
                $+$ & \umlfld{Количество кадров} 
                    &:& \umldtp{Число};\\
                \hdashline
               $+$ & \umlfld{Чересстрочноcть} 
                    &:& \umldtp{\{Да, Нет\}};\\
                \hdashline
                $+$ & \umlfld{Формат} 
                    &:& \umlctp{Формат пикселей};\\
                \hdashline
                $+$ & \umlfld{Есть ссылочные кадры} 
                    &:& \umldtp{\{Да, Нет\}};\\
                \hdashline
                $+$ & \umlfld{Размер группы ссылочных кадров (GOP)} 
                    &:& \umldtp{Число};\\
            \end{tabular}

        }

        \umlclass[
            below right=1em and 0em of Визуальный поток.south west,
        ]{Звуковой поток}{
            \begin{tabular}{cF{7em}cT{6.5em}}
                $+$ & \umlfld{Частота отсчётов}          
                    &:& \umldtp{Число};\\
                \hdashline
                $+$ & \umlfld{Количество отсчётов}   
                    &:& \umldtp{Число};\\
                \hdashline
                $+$ & \umlfld{Количество каналов}
                    &:& \umldtp{Число};\\
                \hdashline
                $+$ & \umlfld{Формат}           
                    &:& \umlctp{Формат отсчётов};\\
                \hdashline
                $+$ & \umlfld{Компоновка}   
                    &:& \umlctp{Формат каналов};\\
            \end{tabular}

        }

        \umlclass[
            left=2em of Звуковой поток,
        ]{Поток субтитров}{}
               
        \umlclass[
            right=3em of Абстрактный поток,
            draw=emphDraw,
            fill=emphClass,
        ]{Пакет}{
            \begin{tabular}{cF{6em}cT{8em}}
                $+$  & \umlfld{Номер}
                    &:& \umldtp{Число};\\
                \hdashline
                $+$ & \umlfld{Время начала} 
                    &:& \umlctp{Точка времени};\\
                \hdashline
                $+$ & \umlfld{Время конца} 
                    &:& \umlctp{Точка времени};\\   
                \hdashline               
                $\#$ & \umlfld{Кадры} 
                     &:& \umlctp{Список кадров};\\
            \end{tabular}
        }

        \umlclass[
            above=1em of Пакет,
        ]{Кодек}{
            \begin{tabular}{cF{4em}cT{8em}}
                $+$ & \umlfld{Название} 
                    &:& \umldtp{Строка};\\
                \hdashline
                $+$ & \umlfld{Кодер} 
                    &:& \umldtp{\{Да, Нет\}};\\
                \hdashline
                $+$ & \umlfld{Декодер} 
                    &:& \umldtp{\{Да, Нет\}};\\
                \hdashline
                $+$ & \umlfld{Опции} 
                    &:& \umldtp{Список опций};\\
            \end{tabular}
        }

        \umlclass[
            right=4em of Кодек,
        ]{Опция кодирования}{
            \begin{tabular}{cF{6em}cT{4em}}
                $+$ & \umlfld{Имя} 
                    &:& \umldtp{Строка};\\
                \hdashline
                $+$ & \umlfld{Значение} 
                    &:& \umldtp{Строка};\\
                \hdashline
            \end{tabular}
        }

        \umlclass[
            below=3em of Пакет,
        ]{Точка времени}{
            \begin{tabular}{cF{6em}cT{6.5em}}
                $+$ & \umlfld{PTS — время отображения} 
                    &:& \umldtp{Число};\\
                \hdashline
                $+$ & \umlfld{DTS — время кодирования} 
                    &:& \umldtp{Число};\\
                \hdashline
                $+$  & \umlfld{Масштаб времени}
                     &:& \umldtp{(Число, Число)};\\
            \end{tabular}
        }

        \umlclass[
            below=2em of Точка времени,
        ]{Тип кадра}{
            \begin{tabular}{cF{9em}cT{3em}}
                $+$ & \umlfld{Название} 
                    &:& \umldtp{Строка};\\
                \hdashline
                $+$ & \umlfld{Количество ссылок} 
                    &:& \umldtp{\{0,1,2\}};\\
            \end{tabular}
        }

        \umlclass[
            below=1em of Тип кадра,
        ]{Формат пикселей}{
            \begin{tabular}{cF{8em}cT{4em}}
                $+$ & \umlfld{Название} 
                    &:& \umldtp{Строка};\\
                \hdashline
                $+$ & \umlfld{Количество цветов} 
                    &:& \umldtp{Число};\\
                \hdashline
                $+$ & \umlfld{Битов на пиксель} 
                    &:& \umldtp{Число};\\
            \end{tabular}
        }

        \umlclass[
            below=2em of Формат пикселей,
        ]{Формат отсчётов}{
            \begin{tabular}{cF{7em}cT{4em}}
                $+$ & \umlfld{Название} 
                    &:& \umldtp{Строка};\\
                \hdashline
                $+$ & \umlfld{Битов на отсчёт} 
                    &:& \umldtp{Число};\\
            \end{tabular}
        }

        \umlclass[
            below=1em of Формат отсчётов,
        ]{Формат каналов}{
            \begin{tabular}{cF{7em}cT{4em}}
                $+$ & \umlfld{Название} 
                    &:& \umldtp{Строка};\\
                \hdashline
                $+$ & \umlfld{Представление} 
                    &:& \umldtp{Строка};\\
            \end{tabular}
        }

        \umlclass[
            type=abstract,
            right=3em of Пакет,
            draw=emphDraw,
            fill=emphClass,
        ]{Абстрактный кадр}{
            \begin{tabular}{cF{3em}cT{9em}}
                $+$ & \umlfld{Номер} 
                    &:& \umldtp{Число};\\
                $+$ & \umlfld{Время} 
                    &:& \umlctp{Точка времени};\\  
                $\#$& \umlfld{Данные} 
                    &:& \umldtp{Байтовый массив} ;\\
            \end{tabular}
        }

        \umlclass[
            below right=2em and 0em of Абстрактный кадр.south west,
        ]{Субтитр}{
            \begin{tabular}{cF{5em}cT{5em}}
                $+$ & \umlfld{Ширина} 
                    &:& \umldtp{Число};\\
                \hdashline
                $+$ & \umlfld{Высота} 
                    &:& \umldtp{Число};\\
                \hdashline
                $+$ & \umlfld{Координата по оси X} 
                    &:& \umldtp{Число};\\
                \hdashline
                $+$ & \umlfld{Координата по оси Y} 
                    &:& \umldtp{Число};\\
                \hdashline
                $+$ & \umlfld{Время начала} 
                    &:& \umlctp{Точка времени};\\
                \hdashline
                $+$ & \umlfld{Время конца} 
                    &:& \umlctp{Точка времени};\\
            \end{tabular}
        }

        \umlclass[
            below right=2em and 0em of Субтитр.south west,
        ]{Визуальный кадр}{
            \begin{tabular}{cF{5em}cT{5em}}
                $+$ & \umlfld{Ширина} 
                    &:& \umldtp{Число};\\
                \hdashline
                $+$ & \umlfld{Высота} 
                    &:& \umldtp{Число};\\
                \hdashline
                $+$ & \umlfld{Тип}    
                    &:& \umlctp{Тип кадра};\\
                \hdashline
                $+$ & \umlfld{Формат} 
                    &:& \umlctp{Формат пикселей} ;\\
            \end{tabular}
        }

        \umlclass[
            below right=2em and 0em of Визуальный кадр.south west,
        ]{Звуковой кадр}{
            \begin{tabular}{cF{5em}cT{5em}}
                $+$ & \umlfld{Частота отсчётов}          
                    &:& \umldtp{Число};\\
                \hdashline
                $+$ & \umlfld{Количество отсчётов}   
                    &:& \umldtp{Число};\\
                \hdashline
                \hdashline
                $+$ & \umlfld{Формат}           
                    &:& \umlctp{Формат отсчётов};\\
                \hdashline
                $+$ & \umlfld{Компоновка}   
                    &:& \umlctp{Формат каналов};\\
            \end{tabular}
        }


        \umlimpl
            {Контейнер}
            {Видео}

        \umlaggreg[
            mult1=1,
            mult2=0..1
        ]
            {Контейнер}
            {Мета}

        \umlcompo[
            mult1=1,
            mult2=1,
            geometry=|-,
            anchors=south west and west
        ]
            {Контейнер}
            {Формат контейнера}

        \umlaggreg[
            arg1=1,
            arg2=0..$\ast$,
            geometry=|-,
            anchors=south east and east
        ]
            {Контейнер}
            {Опция контейнера}

        \umlaggreg[
            mult1=1,
            mult2=0..$\ast$
        ]
            {Формат контейнера}
            {Опция контейнера}

        \umlcompo[
            mult1=1,
            mult2=0..$\ast$
        ]
            {Контейнер}
            {Абстрактный поток}

        \umlcompo[
            mult1=1,
            mult2=0..$\ast$
        ]
            {Абстрактный поток}
            {Пакет}

        \umlaggreg[
            mult1=1,
            mult2=0..$\ast$
        ]
            {Пакет}
            {Абстрактный кадр}

        \umlinherit[
            geometry=-|,
            anchors=west and south west
        ]
            {Визуальный поток}
            {Абстрактный поток}

        \umlinherit[
            geometry=-|,
            anchors=west and south west
        ]
            {Звуковой поток}
            {Абстрактный поток}
    
        \umlinherit[
            geometry=-|,
            anchors=west and south west
        ]
            {Поток субтитров}
            {Абстрактный поток}
        
        \umlinherit[
            geometry=-|,
            anchors=west and south east
        ]
            {Визуальный кадр}
            {Абстрактный кадр}

        \umlinherit[
            geometry=-|,
            anchors=west and south east
        ]
            {Звуковой кадр}
            {Абстрактный кадр}

        \umlinherit[
            geometry=-|,
            anchors=west and south east
        ]
            {Субтитр}
            {Абстрактный кадр}
    
        \umlaggreg[
            mult1=2,
            mult2=1
        ]
            {Абстрактный поток}
            {Точка времени}

        \umlaggreg[
            mult1=2,
            mult2=1
        ]
            {Субтитр}
            {Точка времени}

        \umlaggreg[
            mult1=2,
            mult2=1
        ]
            {Пакет}
            {Точка времени}

        \umlaggreg[
            mult1=1,
            mult2=1
        ]
            {Абстрактный кадр}
            {Точка времени}

        \umlaggreg[
            mult1=1,
            mult2=1
        ]
            {Абстрактный поток}
            {Кодек}

        \umlcompo[
            mult1=1,
            mult2=1
        ]
            {Кодек}
            {Опция кодирования}

        \umlHVHaggreg[
            mult1=1,
            mult2=1,
            arm2=-1em,
            anchors=-30 and west
        ]
            {Визуальный поток}
            {Формат пикселей}

        \umlHVHaggreg[
            mult1=1,
            mult2=1,
            arm2=1em,
            anchors=west and east
        ]
            {Визуальный кадр}
            {Формат пикселей}

        \umlHVHaggreg[
            mult1=1,
            mult2=1,
            arm2=1em,
            anchors=west and east
        ]
            {Визуальный кадр}
            {Тип кадра}

        \umlHVHaggreg[
            mult1=1,
            mult2=1,
            arm2=-1em,
            anchors=east and west
        ]
            {Звуковой поток}
            {Формат отсчётов}

        \umlHVHaggreg[
            mult1=1,
            mult2=1,
            arm2=-1em,
            anchors=east and west
        ]
            {Звуковой поток}
            {Формат каналов}

        \umlHVHaggreg[
            mult1=1,
            mult2=1,
            anchors=west and east
        ]
            {Звуковой кадр}
            {Формат отсчётов}

        \umlHVHaggreg[
            mult1=1,
            mult2=1,
            anchors=west and east
        ]
            {Звуковой кадр}
            {Формат каналов}

    \end{tikzpicture}
\end{document}

